problem:  
Fix one of the Aloff–Wallach 7‑manifolds \(W_{p,q} = \mathrm{SU}(3)/\mathrm{U}(1)_{p,q}\) that is known to admit a **homogeneous Einstein metric** \(g_{E}\) (for example \(W_{1,1}\)).  
Let \(\mathcal{M}^+_{p,q}\) denote the space of smooth Riemannian metrics on \(W_{p,q}\) with **positive sectional curvature**, equipped with the \(C^{k,\alpha}\) topology modulo diffeomorphisms and overall scaling.  
For \(g_0 \in \mathcal{M}^+_{p,q}\), consider the **normalized Ricci flow**
\[
\frac{\partial g}{\partial t} = -2(\mathrm{Ric}(g) - \tfrac{1}{7}r(g) g),
\]
where \(r(g) = \tfrac{1}{7}\int_{W_{p,q}} R(g)\,d\mu_g\) is the average scalar curvature.  

Open question:  
Is the homogeneous Einstein metric \(g_E\) **dynamically stable** under the normalized Ricci flow—that is, does there exist a neighborhood \(\mathcal{U}\subset \mathcal{M}^+_{p,q}\) of \(g_E\) such that for every \(g_0\in\mathcal{U}\), the normalized flow \(g(t)\) exists for all \(t\ge0\) and converges (modulo diffeomorphisms and scaling) to \(g_E\)?  
Equivalently, is \(g_E\) a local attractor in the moduli space of metrics on \(W_{p,q}\) with positive sectional curvature?

Why is it a "good" problem:  
This refined problem is concrete, mathematically sound, and positioned precisely at an open frontier of **Ricci flow stability theory** applied to non‑symmetric, positively curved homogeneous spaces.

1. **Correctly posed analytic question.**  
   It restricts to the Aloff–Wallach spaces for which an Einstein metric is known and formulates stability in well-defined analytic terms—existence, convergence, and neighborhood in the \(C^{k,\alpha}\) topology—avoiding undefined notions like “global attractor.”

2. **Bridges homogeneous geometry and analysis.**  
   The question tests whether the algebraically constructed Einstein metric \(g_E\) behaves as a stable fixed point of a nonlinear parabolic flow, linking the representation theory of \(\mathrm{SU}(3)\) (governing the curvature operator near \(g_E\)) with the spectrum of the Lichnerowicz Laplacian, a deep analytic object.

3. **Beyond symmetry: the first non‑symmetric test case.**  
   Stability of homogeneous Einstein metrics is understood for symmetric spaces, but for non‑symmetric, positively curved homogeneous manifolds—Aloff–Wallach among them—it remains unresolved.  
   Determining linear or nonlinear stability here would push the frontier of Ricci flow dynamics on exotic geometry domains.

4. **Potential geometric consequences.**  
   If \(g_E\) is dynamically stable, small perturbations of the metric (for example those arising from numerical approximation or Cheeger deformation) would asymptotically evolve back to the canonical Einstein metric.  
   Instability, on the other hand, could signal the presence of new Ricci solitons or blow‑up phenomena on compact positively curved spaces, an entirely unexplored phenomenon.

5. **Simplicity with depth.**  
   The statement—*Is the homogeneous Einstein metric on \(W_{p,q}\) stable under normalized Ricci flow?*—is phrased in one sentence, yet its resolution involves sophisticated tools from geometric analysis, representation theory, and curvature computation on homogeneous manifolds.  
   It epitomizes the “narrow entrance, profound depth” character of a truly **good** research‑level problem.